\chapter{Instrukcja uruchomienia projektu}

\section{Wymagania systemowe}

Do uruchomienia projektu potrzebny jest komputer z systemem Windows. Projekt najlepiej uruchomić w Visual Studio 2022. W systemie powinien być zainstalowany .NET SDK oraz składnik \texttt{.NET desktop development}. Nie trzeba instalować osobnego serwera bazy danych, ponieważ projekt korzysta z SQLite.

\section{Sposób konfiguracji środowiska}

Najpierw należy zainstalować Visual Studio 2022. Podczas instalacji trzeba zaznaczyć składnik \texttt{.NET desktop development}. Następnie należy rozpakować projekt i otworzyć plik \texttt{HotelManagementApp.csproj} albo całe rozwiązanie, jeżeli zostało dodane.

Po otwarciu projektu Visual Studio powinno automatycznie przywrócić pakiety NuGet. Jeżeli tego nie zrobi, należy otworzyć terminal w folderze projektu i wpisać:

\begin{lstlisting}[style=textStyle,caption={Przywrócenie pakietów NuGet},label={lst:restore}]
dotnet restore
\end{lstlisting}

\section{Instrukcja uruchomienia aplikacji}

Aby uruchomić projekt od zera, należy wykonać następujące kroki:

\begin{enumerate}
  \item Rozpakować archiwum z projektem.
  \item Otworzyć folder \texttt{HotelManagementApp}.
  \item Uruchomić plik \texttt{HotelManagementApp.csproj} w Visual Studio.
  \item Poczekać, aż zostaną przywrócone pakiety NuGet.
  \item Wybrać projekt \texttt{HotelManagementApp} jako projekt startowy.
  \item Nacisnąć \texttt{F5} albo przycisk \texttt{Start}.
  \item Po uruchomieniu aplikacji można korzystać z zakładek w głównym oknie.
\end{enumerate}

Jeżeli pojawi się problem z pakietami, należy wykonać polecenie \texttt{dotnet restore}. Jeżeli pojawi się problem z bazą danych, można usunąć plik \texttt{hotel.db} i uruchomić projekt ponownie.

\section{Informacje dotyczące bazy danych}

Aplikacja korzysta z SQLite. Nazwa pliku bazy danych to \texttt{hotel.db}. Baza jest tworzona i aktualizowana automatycznie przy uruchomieniu programu przez metodę \texttt{Database.Migrate()}.

Połączenie z bazą znajduje się w pliku \texttt{App.xaml.cs}. W projekcie użyto connection stringa:

\begin{lstlisting}[style=textStyle,caption={Connection string SQLite},label={lst:connectionstring}]
Data Source=hotel.db
\end{lstlisting}

\section{Dane logowania testowego}

W obecnej wersji aplikacja nie posiada logowania, dlatego nie ma danych logowania testowego. Program uruchamia się od razu po włączeniu i użytkownik ma dostęp do wszystkich modułów.

\section{Miejsce na repozytorium GitHub}

W tym miejscu należy wkleić link do repozytorium GitHub z projektem:

\begin{lstlisting}[style=textStyle,caption={Miejsce na link do repozytorium GitHub},label={lst:github}]
https://github.com/Frreezzz/HotelManagementApp/tree/main
\end{lstlisting}
